
homomorphism factors as:
@@M0@@
and the first homomorphism is injective. Since, by VI 4.3, 4.6 and 4.7, @@M1@@ is free over @@M2@@, this implies the injectivity of:
@@M3@@

Thus let @@M4@@. Proceed by induction on the rank of @@M5@@. The case @@M6@@ being trivial, suppose the property proved for @@M7@@. Introduce @@M8@@ and the kernel @@M9@@ of the canonical homomorphism @@M10@@, where @@M11@@ is the structural morphism of @@M12@@. Then @@M13@@ identifies with the scheme of flags of @@M14@@ over @@M15@@, and the induction hypothesis shows that it is enough to prove 3.1 for @@M16@@. Taking VI 1.1 into account, the result to be proved can therefore be stated as follows:

\sgaStableTarget{sga6.e09.cor.3-2}\textbf{Corollary 3.2.} \uline{Let} @@M17@@ \uline{be a noetherian scheme}, @@M18@@ \uline{a locally free} @@M19@@\uline{-Module of rank} @@M20@@, @@M21@@ \uline{the associated projective bundle}, @@M22@@ \uline{the structural morphism of} @@M23@@, \uline{and} @@M24@@ \uline{the class in} @@M25@@ \uline{of the sheaf} @@M26@@. \uline{Then the homomorphism}:
@@M27@@
\uline{defined by}:
@@M28@@
\uline{is an isomorphism}.

We have already seen that this homomorphism is surjective, since the @@M29@@, @@M30@@, form a basis of @@M31@@ over @@M32@@; to see that it is injective, use the homomorphism @@M33@@, defined because @@M34@@ is projective.

